\documentclass[10pt,letterpaper]{article}
\usepackage{cornell-notes}

\title{Clause 08.06.02: The While Statement}
\author{}
\date{}
\setDocTitle{Clause 08.06.02: The While Statement}
\setDocSubtitle{C++ 2024}
\setDocOwner{C++ 2024 Cornell Notes}
\setDocDate{\today}
\setCornellCollection{C++ 2024 Cornell Notes}
\setCornellUnitType{Clause}
\setCornellUnitNumber{08.06.02}
\setCornellUnitTitle{The While Statement}
\hypersetup{pdftitle={Clause 08.06.02: The While Statement},pdfsubject={Cornell notes},pdfauthor={}}

\begin{document}
\maketitle
\begin{CornellOverviewBox}{Overview}
\textbf{Classification:} Core language - Normative\\[2pt]
\textbf{Learning objective:} Understand the rules, terminology, and semantic consequences associated with The while statement.
\end{CornellOverviewBox}\begin{CornellSummaryBox}{Summary}
{\sffamily\bfseries\color{Accent} Summary}\par\vspace{3pt}Section 8.6.2 focuses on The while statement. Use the listed grammar productions together with the semantic constraints and disambiguation rules referenced by the section. When applying it, preserve the distinction between mandatory requirements, recommendations, permissions, and behavior deliberately left undefined, unspecified, or implementation-defined.
\end{CornellSummaryBox}

\end{document}
